\section{A common posterior bound over dyadic Gaussian scales}
\label{app:weighted-code}

The Gaussian-weighted code can be tuned to a fixed upper scale.
Tuning the variance sample count to each fixed upper scale keeps the
Gaussian row cost below the finite-parameter cost. A finite sequence
of dyadic scale intervals suffices for this choice. The resulting bound is simultaneous in learner
width and Gaussian scale.

Use the posterior resources in
\eqref{eq:factor-posterior-budgets}, with the product-cylinder cost
$D_*$, so that the retained posterior has density at most
$\exp(NK/t^2)$ relative to its raw Gaussian law. In particular
$\beta=\varepsilon/4$, $\alpha=\varepsilon/12$, and
$e=\alpha r^2/(1000L)<1$.
The rank cap $q$, message charge $D_0$, and the pilot, correction, and
readout counts do not depend on the scale ceiling. Only the variance
count and its consequent scalar mesh need to change.

Put
\begin{equation}\label{eq:weighted-dyadic-scales}
 J_*=\lfloor\log_2(1/e)\rfloor,\qquad
 v_j=2^{-j}\quad(0\le j\le J_*),\qquad B_*=J_*+1.
\end{equation}
Every $u\in(e,1]$ has an index $j$ for which
$v_j/2<u\le v_j$, and every listed $v_j$ lies in $[e,1]$.
For each $j$, use the weighted code with ceiling $t_0=\sqrt{v_j}$.
Write $k_j$ for its variance count, $J_{b,j}$ for its finite-code budget
\eqref{eq:weighted-finite-budget}, $p_j$ for its Gaussian dimension
\eqref{eq:weighted-gaussian-dimension}, and $\Pi_{j,u}$ for its
code reference prior. All these objects are fixed before observing data.
For the cold range $0<u\le e$, let $J_c$ be the ordinary finite-code
budget (EM25) at ceiling $t_0=\sqrt e$.

The information threshold and the common charge are
\begin{equation}\label{eq:weighted-dyadic-budget}
 \begin{split}
 \Delta&=\max\left\{K,\ 2LD_0,\
                  \max_{0\le j\le J_*}\frac{2Kk_j}{v_j}\right\},\\
 \mathcal B&=\max\left\{\Delta+1,\ J_c,\
       \max_{0\le j\le J_*}
       \left[J_{b,j}+\frac{4Lk_jK}{v_j}+1\right]\right\}.
 \end{split}
\end{equation}
Let $M,R_\tau$ be the Gaussian reference-family sizes in
Lemma~\ref{lem:si-grid}, with its parameter $J=\Delta$.
For each dyadic interval also define
\begin{equation}\label{eq:weighted-dyadic-grid}
 R_j=1+\left\lfloor\frac{p_j}{2}\log2\right\rfloor,\qquad
 \tau_{j,\ell}=\sqrt{v_j}\exp(-\ell/p_j)
          \quad(0\le\ell<R_j),
\end{equation}
and put
\[
 R_{\max}=\max_j R_j,\qquad
 \Lambda=\log\frac{3[MR_\tau+B_*R_{\max}+1]}{\delta}.
\]

\begin{theorem}[Weighted-code learning at every mean-field-or-stronger scale]
\label{thm:weighted-all-scales}
Under the model and Boolean-teacher conventions of the paper, the
sample condition
\begin{equation}\label{eq:weighted-sample-budget}
 n\ge\frac4\varepsilon[\mathcal B+\Lambda]
\end{equation}
implies, with probability at least $1-\delta$ over the iid training
sample,
\[
 R(Q_{S,N,t,c})\le\varepsilon
 \qquad\text{simultaneously for every }N\ge C,
                      \quad0<t\le1.
\]
The same statement holds after adding the label-independent prior
restriction $|z_w(x)|>1/2$ at every Boolean input.
\end{theorem}
\begin{proof}
We use one sample-independent reference prior on a tagged predictor
space. Give weight $1/3$ to each of the following parts:
\begin{enumerate}
\item the finite Gaussian family for $NK/t^2\le\Delta$;
\item the uniform law on the ordinary cold finite class;
\item a uniform choice of $j$, then a uniform choice of
$\ell\in\{0,\ldots,R_j-1\}$, followed by the weighted code reference
law $\Pi_{j,\tau_{j,\ell}^2}$.
\end{enumerate}
Use the test-margin indicator loss on the raw Gaussian part and the
ramp loss on both comparison-code parts. These form one bounded loss
on the tagged space. Apply Lemma~\ref{lem:si-risk} once.

Write $u=t^2$ and $b=NK/u$. If $b\le\Delta$, the finite Gaussian
family applies to the actual posterior. Its empirical indicator loss
is zero. Including the mixture weight and confidence term, its charge
is at most $\mathcal B+\Lambda$, so its risk is at most
$2(\mathcal B+\Lambda)/n$.

Suppose $b>\Delta$. Then
$N/u>\max\{1,2LD_0/K\}$, as required by both compression arguments.
The readout-tail estimate retains mass at least $1-\beta$ on the common
robustness radius $r$, and its normalized retained law satisfies the
same density bound with the specified $K$.
If $u\le e$, apply the ordinary code at ceiling $\sqrt e$. Its
relative-entropy charge is at most $J_c$ before mixture weights.

If $u>e$, choose $j$ with $v_j/2<u\le v_j$. From
\eqref{eq:weighted-dyadic-budget},
\[
 N=bu/K>\frac{2k_j u}{v_j}>k_j.
\]
Thus sampling the variance group without replacement is permitted.
Theorem~\ref{thm:weighted-code-kl} gives a law on the $j$th code space
with relative entropy at most
\[
 J_{b,j}+2Lk_jK/u
 \le J_{b,j}+4Lk_jK/v_j.
\]
Choose $\ell=\lfloor p_j\log(\sqrt{v_j}/t)\rfloor$.
Then $0\le\ell<R_j$, $\tau_{j,\ell}\ge t$, and
$p_j\log(\tau_{j,\ell}/t)<1$ (with a non-strict bound also sufficient).
For centered isotropic Gaussians in $p_j$ dimensions,
\[
 \frac{d\mathcal N(0,t^2I)}
      {d\mathcal N(0,\tau_{j,\ell}^2I)}\le\exp(1).
\]
The quantizer and remaining finite-code reference law are fixed within
this dyadic interval. Pushing the domination through the quantizer
therefore changes the relative-entropy charge by at most one. The
mixture and confidence costs are bounded by $\Lambda$.

In either compression branch, the common pointwise coupling gives
empirical ramp loss at most $\alpha$ and
$R(Q_{S,N,t,c})\le\beta+\alpha+R_\psi(\Gamma)$.
Lemma~\ref{lem:si-risk} consequently bounds the desired risk by
\[
 \beta+3\alpha+\frac2n[\mathcal B+\Lambda]
 \le\varepsilon.
\]
All comparisons use the same sample event. Optional global-margin
conditioning retains the same target cylinder, readout-tail bound,
and canonical compression argument, so the proof is unchanged.
\end{proof}

\begin{corollary}[Quartic width and ninth-power depth]
\label{cor:weighted-powers}
Fix $d,c,\varepsilon,\delta$ and put
$F_0=1+\log^2(2+C)$ and $G=\log(2+CL)$.
There is a constant $A_0$ depending only on these fixed quantities such
that
\[
 n\ge A_0 C^4L^9F_0^4G^2
\]
suffices in Theorem~\ref{thm:weighted-all-scales}.
\end{corollary}
\begin{proof}
At fixed input dimension, suppressing only logarithms in the explicit
resources, the finite part of every dyadic code has
\[
 J_{b,j}=\widetilde O(LK^2/e^2).
\]
The finite Gaussian row dimension affects the scale-grid charge only
logarithmically. Since $v_j\in[e,1]$, the explicit variance count obeys
\[
 \frac{k_j}{v_j}
 =\widetilde O\left(\frac1{v_j}+\frac{v_j}{e^2}\right)
 =\widetilde O(e^{-2}).
\]
It follows that
\[
 \Delta=\widetilde O(LK^2/e^2),\qquad
 \mathcal B=\widetilde O(LK^2/e^2).
\]
For the cold code the scale term in (EM26) is only $LK/e$, so it is
also absorbed by this bound. These resource statements keep input
dimension fixed; the exact formulas above retain its full dependence.

For the target-cylinder budgets,
\[
 K=O(CLF_0),\quad e^{-1}=O(CL^3F_0),\quad
 q=O(C^2L^4F_0^2),\quad m_b=O(qG).
\]
The finite parameter count is $O(Lq^2G)$, and its scalar mesh costs
$O(G)$ per parameter. Thus
$J_{b,j}=O(C^4L^9F_0^4G^2)$ uniformly in $j$.
The weighted-row charge and raw cutoff are smaller:
\[
 \max_j Lk_jK/v_j=O(C^3L^8F_0^3G),\qquad
 LD_0=O(C^4L^9F_0^4G).
\]
The ordinary cold code has the same leading finite-parameter order,
and all remaining scale-family charges are logarithmic. Substituting
these estimates in \eqref{eq:weighted-sample-budget} proves the claim.
\end{proof}
